\scanpage{51}
(ii) Ce raisonnement vaut aussi sans supposer $T/S$
affine, si $g$ est \uline{étale}, \deleted{\unsure{$f_\ast$}} Ainsi, $f_\ast(g)$ sera
\deleted{\unsure{ét.}}simple. Pour prouver qu'il est étale, il suffit
de prouver que \unsure{les|ses} \uline{fibres sont} \deleted{\unsure{de}} \inserted{\uline{finies}
(\unsure{les fibres|ensuite})} \deleted{de dim $0$ (\unsure{à fibres} \illegible{1})}. On est ramené
au cas où $S=\operatorname{Spec}(k)$, $k$ un corps, et
% slanted
\margin{\unsure{Donner un|Donne un|Dans un} \uline{corollaire} en fait}
\diagram{\begin{tikzcd}[ampersand replacement=\&]
\& X\arrow[r]\&Y\\
k\&T\&
\end{tikzcd}}
\diagram{\begin{tikzcd}[ampersand replacement=\&]
\& f_\ast(X)\arrow[d,very thick]\\
\operatorname{Spec}(k')\arrow[r]\& f_\ast(Y)
\end{tikzcd}}
à prouver que si $k'$ est une
clôture algébrique de $k$, alors il
n'y a qu'un nb fini de $k'$-pts de $f_\ast(X)$ au dessus d'un
$k'$-point de $f_\ast(Y)$, i.e. \deleted{\unsure{que}} il n'y a qu'un nb fini
de sections de $X'/T'$ qui
relèvent une section donnée de $Y'/T'$.
[les $'$ \unsure{dénotent|dit} ext. des scalaires de $k\unsure{:|\text{ à }}k'$].
Or $T'$ sera noethérien [de type fini sur $k'$]
donc n'a qu'un nb fini de comp. connexes.
On peut supposer $T'$ connexe. Alors le
relèvement est connu quand on
connaît l'image d'un pt. On trouve
donc

\uline{Corollaire} $1$ Soit $\unsure{\nu|v}$ le maximum des nbs
\deleted{\unsure{comp.|rangs}} de composantes connexes absolues
des $T_s$ ($s\in S$), et soit \deleted{\unsure{de}} $=$ $n$ le
maximum analogue pour $X/Y$. Alors le maximum
% archive sheet number: 50
